% ══════════════════════════════════════════════════════
%  contenu/exos-resolus.tex  —  Exercices résolus
% ══════════════════════════════════════════════════════
\clearpage
\PartieHeader{Exercices r\'esolus}

\begin{ExoResolu}
\textbf{Probl\`eme de m\'elange.}
Un r\'eservoir contient initialement 500 litres d'eau pure. Il est aliment\'e par un
conduit qui fournit de l'eau \`a 4~L/min avec une concentration de sel de 100~g/L.
Un autre conduit laisse l'eau sal\'ee s'\'echapper \`a 4~L/min.
\begin{enumerate}
  \item D\'eterminer $q(t)$, la quantit\'e de sel pr\'esente dans le r\'eservoir en fonction du temps.
  \item Donner la quantit\'e de sel apr\`es 10 minutes et apr\`es 1 heure.
\end{enumerate}
\end{ExoResolu}

\begin{Correction}
Soit $q(t)$ la quantit\'e de sel (en kg) dans le r\'eservoir au temps $t$ (en minutes).

\textbf{Mod\'elisation :} $\dfrac{dq}{dt}=\text{input}-\text{output}$.

\textit{Input} $=4\,\text{L/min}\times 100\,\text{g/L}=0{,}4\,\text{kg/min}$.

\textit{Output} $=\dfrac{q(t)}{500}\,\text{kg/L}\times 4\,\text{L/min}=\dfrac{q(t)}{125}\,\text{kg/min}$.

L'\'equation diff\'erentielle est :
$q'+\dfrac{q}{125}=0{,}4$ avec $q(0)=0$.

\textbf{Solution g\'en\'erale :} $q(t)=50+Ce^{-t/125}$.

\textbf{Condition initiale} $q(0)=0$ : $C=-50$, donc :
\formula{q(t)=50\!\left(1-e^{-t/125}\right)}

Apr\`es 10 min : $q(10)=50(1-e^{-2/25})\approx 3{,}84$ kg.
Apr\`es 1 h : $q(60)\approx 19{,}06$ kg.\hfill$\square$
\end{Correction}

\begin{ExoResolu}
\textbf{Charge d'un condensateur.}
Un condensateur $C=4\times 10^{-4}$ F en s\'erie avec $R=88\,\Omega$ et $E=6$ V.
La tension $u_c$ v\'erifie $E=RC\dfrac{du_c}{dt}+u_c$, avec $u_c(0)=0$.
\begin{enumerate}
  \item \'Ecrire l'\'equation sous la forme $y'=ay+b$.
  \item R\'esoudre avec la condition initiale.
  \item Donner $u_c$ apr\`es $t=100$ ms.
\end{enumerate}
\end{ExoResolu}

\begin{Correction}
\textbf{1)} $y'=-\dfrac{1}{RC}y+\dfrac{E}{RC}$, avec $a=-\dfrac{1}{RC}$, $b=\dfrac{E}{RC}$.

\textbf{2)} Solution g\'en\'erale : $y=ke^{at}+E$. Condition $u_c(0)=0$ : $k=-E$.
\formula{u_c(t)=E\!\left(1-e^{-t/(RC)}\right)=6\!\left(1-e^{-t/0{,}0352}\right)}

\textbf{3)} $u_c(0{,}1)=6(1-e^{-125/44})\approx 5{,}65$ V.\hfill$\square$
\end{Correction}

\begin{ExoResolu}
\begin{enumerate}
  \item R\'esoudre $y'-y=0$ \quad $(1)$.
  \item Soit $f$ d\'erivable sur $\R$ avec $f'(x)-f(x)=x$ \quad $(2)$.
    \begin{enumerate}[label=\alph*)]
      \item D\'eterminer $a,b$ pour que $f_0(x)=ax+b$ soit solution de $(2)$.
      \item D\'eduire toutes les solutions de $(2)$ (poser $y=f-f_0$).
      \item D\'eterminer l'unique solution $f_1$ v\'erifiant $f_1(0)=0$.
    \end{enumerate}
\end{enumerate}
\end{ExoResolu}

\begin{Correction}
\textbf{1)} Solution g\'en\'erale de $y'=y$ : $y=\lambda e^x$, $\lambda\in\R$.

\textbf{2a)} $f_0'-f_0=a-(ax+b)=-ax+(a-b)=x$.

Identification : $\begin{cases}-a=1\\a-b=0\end{cases} \Rightarrow a=-1,\;b=-1$.
Donc $f_0(x)=-x-1$.

\textbf{2b)} $y=f-f_0$ v\'erifie $y'-y=0$, donc $f-f_0=\lambda e^x$.
\formula{f(x)=-x-1+\lambda e^x, \quad \lambda\in\R}

\textbf{2c)} $f_1(0)=0$ : $-1+\lambda=0 \Rightarrow \lambda=1$, donc $f_1(x)=e^x-x-1$.\hfill$\square$
\end{Correction}

\begin{ExoResolu}
R\'esoudre $2y''-5y'+2y=0$, puis d\'eterminer la solution
v\'erifiant $y(0)=1$ et $y'(0)=0$.
\end{ExoResolu}

\begin{Correction}
\'Equation caract\'eristique : $2r^2-5r+2=0$, $\Delta=9$.
$r_1=2$, $r_2=\dfrac{1}{2}$.

Solution g\'en\'erale : $y=\lambda e^{2x}+\mu e^{x/2}$, $(\lambda,\mu)\in\R^2$.

Conditions : $\begin{cases}\lambda+\mu=1\\2\lambda+\frac{\mu}{2}=0\end{cases}$
$\Rightarrow \lambda=-\dfrac{1}{3}$, $\mu=\dfrac{4}{3}$.
\formula{y(x)=-\frac{1}{3}e^{2x}+\frac{4}{3}e^{x/2}}
\hfill$\square$
\end{Correction}

\begin{ExoResolu}
R\'esoudre $(E_m):y''-2y'+(1-m)y=0$ selon les valeurs du r\'eel $m$.
\end{ExoResolu}

\begin{Correction}
\'Equation caract\'eristique : $r^2-2r+(1-m)=0$, $\Delta=4m$.

\begin{center}
\renewcommand{\arraystretch}{1.7}
\begin{tabular}{|c|c|c|}
\hline
\y \textbf{Cas} & \y \textbf{Racines} & \y \textbf{Solution g\'en\'erale} \\
\hline
$m=0$ & racine double $r_0=1$ & $y=(\lambda+\mu x)e^x$ \\
\hline
$m>0$ & $r_1=1+\sqrt{m}$, $r_2=1-\sqrt{m}$ & $y=\lambda e^{(1+\sqrt{m})x}+\mu e^{(1-\sqrt{m})x}$ \\
\hline
$m<0$ & $r=1\pm i\sqrt{-m}$ & $y=e^x[\lambda\cos(\sqrt{-m}\,x)+\mu\sin(\sqrt{-m}\,x)]$ \\
\hline
\end{tabular}
\end{center}

Dans tous les cas $(\lambda,\mu)\in\R^2$.\hfill$\square$
\end{Correction}
